\documentclass{ltjsbook}
\usepackage{docmute} 
\usepackage{../../myStyle} 

\begin{document}
% \chapter{テスト理論と因子分析}
\chapter{Test Theory and Factor Analysis}
\label{x03_testTheory}

% \section{古典的テスト理論}
\section{Classical Test Theory}
% 前回の信頼性についての議論のなかで，古典的テスト理論について触れました。古典的テスト理論のモデルは$X=t+e$で表されます。すなわち，テストの点数$X$は真のスコア$t$と誤差$e$に分割できるというものです。
In the discussion about reliability last time, we touched on classical test theory. The model of classical test theory is expressed as $X=t+e$. That is to say, the test score $X$ can be split into the true score $t$ and the error $e$.
% 非常に単純なモデルですが，測定したものには誤差がついているという考え方，言い換えると目に見えるものだけが真実ではないという考え方がしめされています。この考え方は，ソフトサイエンスの領域においては重要なことです。
It's a very simple model, but it encapsulates the idea that what is measured comes with errors, or to put it another way, not everything visible is the truth. This idea is important in the field of soft science.

% またこの古典的テスト理論から，いくつかの重要な考えを読み取ることができます。ひとつは誤差についての考え方です。このモデルを$X_j = t_j + e_j$のように，ある個人の変化しない属性について，$j$回測定したとします。この時，測定の平均は次のように計算できます\footnote{この式は，ある測定を多くの個人$i$に対して行ったものとして，$X_i = t_i + e_i$と考えることもできますが，添字が異なるだけで式の展開に違いはありません。}。
From this classical test theory, we can derive some important ideas. One is the concept of errors. Suppose we measure an individual's unchanging attribute in this model as $X_j = t_j + e_j$, where $j$ is the number of measurements. In this case, the average of the measurements can be calculated as follows\footnote{This equation can also be interpreted as $X_i = t_i + e_i$ where the measurement is performed on many individuals $i$, but there is no difference in the expansion of the equation other than the difference in subscripts.}.

\begin{align*}
     & \bar{X}  =  \frac{1}{n}\sum_{j=1}^n X_j                           &             & \mbox{} \\
     & =  \frac{1}{n}\sum_{i=j}^n (t_j + e_j)                            & \mbox{定義より}           \\
     & =  \frac{1}{n} \sum_{j=1}^n  t_j +  \frac{1}{n} \sum_{j=1}^n  e_j & \mbox{分配して}           \\
     & =  \bar{t} + \bar{e}                                              & \mbox{}               \\
\end{align*}


% さて古典的テスト理論では，誤差に関して次のことが仮定されます。
In classical test theory, the following assumptions are made about errors.
\begin{itemize}
	\item 誤差の平均はゼロ。つまり誤差が出現するときは，「正に偏る」「負に偏る」といった一貫した傾向がないと考える。
	\item 真のスコアと誤差との相関はゼロ。つまり誤差は真のスコアに関係なく影響してくるもので，真のスコアと共変動するようであれば偶然によるものとはいえない。
	\item 異なる測定誤差間の相関はゼロ。誤差同士がなにか意味のある変動をしているのであれば，それはもう偶然によるものとはいえない。
\end{itemize}


% これらはいずれも，誤差が「偶然によって現れる影響で，制御不可能なもの」という考え方からは自然な仮定だといえるでしょう。より詳しくいえば，誤差は測定に応じて毎回一定の傾向で生じる\keyterm{系統誤差}{けいとうごさ}{systematic error}と，全く傾向のつかめない\keyterm{偶然誤差}{ぐうぜんごさ}{random error}に分けて考えられますが，系統誤差は測定に際して工夫して取り除くべき問題であり，ここでは全くの偶然による誤差についての議論だからです。
All of these can be said to be natural assumptions from the perspective that errors are "effects that appear by chance and are uncontrollable." To elaborate more specifically, errors can be considered as either systematic errors, which occur consistently according to each measurement, or random errors, which are entirely unpredictable. Systematic errors are issues that should be eliminated with careful consideration during measurement, which is why the discussion here is entirely about errors caused by chance.

% さて誤差の平均がゼロ，つまり$\bar{e}=0$ですから，$\bar{X}=\bar{t}$となって，いつかは誤差がなくなって真のスコアを得ることができるようになる，ということが示されます。
Well, since the average error is zero, i.e., $\bar{e}=0$, it follows that $\bar{X}=\bar{t}$, demonstrating that eventually the error will be eliminated and the true score can be obtained.


% また平均は$0$ですが分散はゼロではありません\footnote{ガウスの考えた誤差論から，誤差は確率\keyterm{正規分布}{せいきぶんぷ}{Gaussian Curve}に従うと考えられます。}。このテストの分散を考えると，次のようなことがわかります。
Although the mean is $0$, the variance is not zero\footnote{According to Gaussian error theory, errors are considered to follow the probability \keyterm{Normal Distribution}{Normal Distribution}{Gaussian Curve}}. Considering the variance of this test, we understand the following.
\begin{align*}
	 & Var(X)  =  \frac{1}{n}\sum_{i=j}^n (X_j - \bar{X})^2                                                                                              & \mbox{定義より}         \\
	 & = \frac{1}{n}\sum_{j=1}^n  ((t_j + e_j) - (\bar{t} + \bar{e}))^2                                                                                  & \mbox{Xをテスト理論のモデルに} \\
	 & = \frac{1}{n}\sum_{j=1}^n (  (t_j - \bar{t}) + (e_j - \bar{e}) )^2                                                                                & \mbox{同じ記号でまとめる}    \\
	 & = \frac{1}{n}\sum_{j=1}^n (  (t_j - \bar{t})^2 + 2(t_j - \bar{t})(e_j - \bar{e}) + (e_j - \bar{e})^2  )                                           & \mbox{展開する}         \\
	 & = \frac{1}{n}\sum_{i=1}^n (t_j - \bar{t})^2 +\frac{1}{n}\sum_{j=1}^n 2(t_j - \bar{t})(e_j - \bar{e})  + \frac{1}{n}\sum_{j=1}^n (e_j - \bar{e})^2 & \mbox{$\sum$を分配}    \\
	 & =  Var(t) + 2Cov(te) + Var(e)
\end{align*}

% ここで$Cov$とは共分散を表しています。第二項の$2Cov(te)$は真のスコアと誤差との共分散(を2倍したもの)を表していることになりますが，共分散(相関)がそもそも線形関係を表す指標であったことを思い出してください。相関係数は共分散を標準化したものだったわけですが，そういう意味ではこの$Cov(te)$というのは真のスコアと誤差との相関関係を表しているようなものです。さて，ここでも誤差の仮定から，相関はゼロです。すなわち，誤差とはどのような傾向もなく出現するもの，という考えられているのです。どのような傾向もないわけですから，当然なにかと相関関係にあるはずがない，すなわち$Cov(te)=0$であるとなります。
Here, $Cov$ represents covariance. The second term $2Cov(te)$ represents the covariance (doubled) between the true score and the error, but remember that covariance (or correlation) is an indicator of a linear relationship in the first place. Correlation coefficients were standardized covariances, but in that sense, this $Cov(te)$ represents a correlation between the true score and the error. Now, here too, from the assumption of error, the correlation is zero. In other words, it is thought that errors occur without any trend. Since there is no trend, it naturally should not be correlated with anything, that is $Cov(te)=0$.

% これを踏まえると，$Var(X)=Var(t)+Var(e)$のように，テストの分散が真のスコアの分散と誤差の分散に完全に分割されました。ここから，信頼性の定義は$\frac{Var(t)}{Var(X)}=\frac{Var(t)}{Var(t)+Var(e)}$と表すことができるのですね。
Taking this into account, the variance of the test, $Var(X)$, is completely divided into the variance of the true score $Var(t)$, and the variance of the error, $Var(e)$. From here, we can represent the definition of reliability as $\frac{Var(t)}{Var(X)}=\frac{Var(t)}{Var(t)+Var(e)}$.
% 言葉で言えば，\keytermL{信頼性}{しんらいせい}の定義は「全分散中にしめる真のスコアの分散の割合」ということになります。
In words, the definition of \keytermL{reliability}{しんらいせい} is "the proportion of the true score variance in the total variance."

% ここまでは古典的テスト理論から示されることであり，これまでの復習になります。ここから，このテスト理論をより展開させていくことを考えましょう。
Up to this point, what has been shown is from classical test theory, and it serves as a review. From here on, let's consider developing this test theory further.

% \section{因子分析モデル}
\section{Factor Analysis Model}
% \subsection{単因子モデル}
\subsection{Unifactorial Model}
% 今からお話しするのは，\keyterm{因子分析}{いんしぶんせき}{Factor Analysis}というモデルです。因子分析モデルは古典的テスト理論の拡張であり，もっとも簡単な1因子モデルは次のように表されます。
The model I'm going to talk about now is Factor Analysis. The factor analysis model is an extension of classical test theory, and the simplest one-factor model is represented as follows.

\[ z_{ij} = a_j f_i + e_{ij} \]

ここで$Z_{ij}$は個人$i$の項目$j$に対する反応を標準得点で表したもの\footnote{\keyterm{標準得点}{ひょうじゅんとくてん}{Standard Score}とは，素点$X_j$を$Z_j = \frac{X_j - \bar{X}}{\sigma_j}$と変換したものであることを思い出してください。標準化されたスコアは平均が0,分散が1になりますので，単位の異なるもの同士であっても標準得点を使うと相互に比較可能になるのでした。}，$a_j$は項目$j$の\keyterm{因子負荷量}{いんしふかりょう}{factor loading}，$f_i$は個人$i$の\keyterm{因子得点}{いんしとくてん}{factor score}，$e_{ij}$は個人$i$と項目$j$の組み合わさった時に生じた誤差と呼ばれます。

\keyterm{因子負荷量}{いんしふかりょう}{factor loading}とは，因子というこのテストで測定したい特性と，項目との関係の強さを表しているものです。
\keyterm{因子得点}{いんしとくてん}{factor score}とは，因子というこのテストで測定したい特性と，個人との関係の強さを表しているもので，その人のスコアだということができます。

記号についている添字に注目してください。添字$i$は個人を，添字$j$は項目を表していますが，因子負荷量は$a_j$と表されています。つまり項目によって変わる変量です。因子得点は$f_i$と表されています。つまり人によって変わる変量です。古典的テスト理論をこの添字を使って表現するならば，$X_{ij} = t_i + e_{ij}$となりますが，これと比べてみると$t_i$が$a_jf_i$に変わったのが因子分析モデルだということになります。$t_i$は個人についての真のスコアなのですが，古典的テスト理論の場合，テストの点数は個人のこの能力だけを反映していると考えられていたことになります。もしテストの問題が難しすぎて，まったく答えることができなければ，その人の能力はゼロということになるわけです。しかし中には悪い問題というのもあって，たとえば小学生に高校で習う知識が必要な問題を解かせるような問題があれば，誰だって解けないかもしれません。解けない問題を出して「学力が低い」と結論づけるのはやや暴力的ですらありますね。このように，古典的テスト理論は項目の良し悪しといったものが評価できないモデルだったのです。

因子分析モデルはこれを改良し，$a_j f_i$としました。すなわち，ある項目に対する反応$z_{ij}$は，その項目が測定したい特徴を十分に反映しているかどうか($a_j$)と，その人がその特徴を有しているかどうか($f_i$)の両方が成立している必要があるわけです。掛け算ですから，一方がゼロであれば結果もゼロになります。すなわち測定したい特徴を反映していない項目($a_j=0$)であれば，どれほどそれについての能力($f_i$)が高くても反応できないのです。たとえば美的センスに非常に秀でた人がいても，数学のテストでその能力を反映させることはできませんよね。これは数学のテストというのが数学力を測定するものであって，美的センスを測定するものではないからです。

因子分析は知能検査や性格検査の文脈から生まれてきたものです。心理学において「知能」とは，何にでも応用可能な一般的な知能というのがあるのか，あるいは語彙力や計算力といった複数の個別の能力があるのか，という議論がありました。知能検査としていろいろなものが考えられますが，それらがきちんと当該能力を測定する検査法だったかどうかもわからないわけで，因子分析モデルはそこを評価できるようにした，とも言えます。性格検査についても同様で，特性論的に考えるならば人間の性質というのは複数のもの，たとえば外向性，神経症傾向，開放性，協調性，勤勉性\footnote{\textcite{Oshio2020}のビッグファイブについての解説に基づいています。}などがあり，ある性格検査の項目は協調性を測定するのには向いているけれども，神経症傾向を測定するには向いていないということがあるわけです。このように，心理学と因子分析モデルは深い関係があります。

\subsection{多因子モデル}
さて先ほどは一因子，あるいは単因子ともいいますが，測定したい特徴が1つだけのモデルでした。学力テストなどは一因子で問題ありません。国語のテストは国語の能力を，数学のテストは数学の能力を測定していれば良いのであって，数学のテストを解くのに美的センス(真美的能力)が必要というのは，むしろ困った状況ですらありますね。しかし，知能検査や性格検査の場合はそうではありません。ある行動傾向，ある形容詞，ある検査がたった1つの能力・性質・心理的要因だけを反映しているとは限りません。たとえば人に優しく振る舞うといっても，その背後に外向性があるのか，あるいはそうすると自分がよく見られるからという利己的な性格があるのか等々が考えられます。1つの項目に複数の要素(因子)が複合的に影響していることを考えるべきです。そこで因子の数が1つではなく，複数ある多因子モデルを考えることにします。多因子モデルは次のように表現されます。

\begin{equation}
	z_{ij} = a_{j1}f_{i1} + a_{j2}f_{i2}+a_{j3}f_{i3} + \cdots + a_{jm}f_{im} + d_j u_{ij}
	\label{factor}
\end{equation}

記号の意味は単因子モデルと基本的には同じです。$z_{ij}$は個人$i$の項目$j$に対する反応を標準得点で表したものであり，$a_{jm}$は第m因子の\keytermL{因子負荷量}{いんしふかりょう}，$f_{im}$は第$m$因子の\keytermL{因子得点}{いんしとくてん}を表しています。因子の数が複数あるモデルですから，$a_{j1},a_{j2},\cdots, a_{jm}$とか$f_{i1},f_{i2},\cdots, f_{im}$のように因子の番号と項目・個人の添字の組み合わせになっていることを確認してください。最後の$e_{ij}$が$d_j u_{ij}$となっていますが，これは誤差についても他の因子と形を同じくし，項目に依存するものとそれ以外に区別しているだけです。

さて，ここでは因子を$m$個あるとしています。この因子はどの項目にも共通して働くので，\keyterm{共通因子}{きょうつういんし}{common factor}と呼ばれます。共通因子がいくつあるかは事前にわかりませんが，一般的にその数は数個〜十数個になります。性格心理学は長い研究の中で，性格を表す言葉に共通する因子はおおよそ5つぐらいであろう，という答えを得るに至りましたが，それ以外の領域では領域ごとの見解があるでしょう。知能が何種類の因子に分かれるのか，あるいはとある心の状態がどういう構造をしているのか，というのは心理学的にみても十分興味のある考え方です。もちろん因子分析によって得られる因子が，人間の潜在的な知能や概念に直接対応しているとは言えないのですが\footnote{むしろテスト項目や調査票などに対する反応パターンが因子として出てくるだけで，性格や知能が数次元あるというより，我々は性格や知能を数次元で捉えることしかできない，という言い方の方が正しいでしょう。}，それでも因子がどのような構造(しくみ)をしているのかについての一定の情報を与えてくれます。多因子モデルが心理学一般で広まったのは，こうした心の「構造」に注目する学問との相性が良かったということでしょう。

\section{因子分析の定理}
\subsection{因子分析モデルの展開}
因子分析モデルも古典的テスト理論のように，式の展開から何が見えてくるか考えてみましょう\footnote{以下このセクションは\textcite{kosugi2018}の原稿を再構成したものです。}。

左辺の$z_{ij}$は観測されたデータから算出できるものですが，右辺の因子負荷量，因子得点はいずれも未知数です。データに対して未知数が多すぎるようで，これではどのようにして答えを見つけ出せば良いのかわからないかもしれません。たとえばある人のある項目に対する標準得点が0.12であるとして，それが$0.4 \times 0.3$で得られるのか，$0.2 \times 0.6$で得られるのか，はたまた他の数値の組み合わせで得られるのか，を解く数学的技術は存在しません。この方程式はこのままでは解けないのです。

そこで，この未知数だらけの方程式を解くために，因子について以下のような条件を置きます。
\begin{itemize}
	\item 共通因子の因子得点，独自因子の因子得点は，標準化されている。すなわち，いずれの因子得点も平均点は0であり，分散は1である。
	\item 独自因子は共通因子，他の独自因子と相関しない。
\end{itemize}
この他に，状況に応じて因子同士の間に相関を仮定します。
\begin{itemize}
	\item 共通因子同士の相関を認めないのを「直交因子モデル」，認めるのを「斜交因子モデル」と呼ぶ。
\end{itemize}

このような仮定を置いたら問題が解けるようになるのでしょうか？ 実はこの問題を解く鍵は，多変量データであればなんとかなるのです！

2つの変数，$j$と$k$の標準得点から，
\begin{equation}
	r_{jk}=\frac{1}{N} \sum_{i=1}^{N} z_{ij}z_{ik}
\end{equation}

のように，相関係数が算出されることを思い出してください。先ほどの因子分析の基本代数式(式\ref{factor})をこの式に代入してみましょう。
\begin{equation}
	\begin{array}{ccl}
		r_{jk} & = & \displaystyle{\frac{1}{N} \sum_{i=1}^{N} z_{ij}z_{ik}}                                                                                                                  \\
		       & = & \displaystyle{\frac{1}{N} \sum_{i=1}^{N}}(a_{j1}f_{i1}+a_{j2}f_{i2}+ \ldots \ +a_{jm}f_{im}+d_{j}u_{ij})(a_{k1}f_{i1}+a_{k2}f_{i2}+ \ldots \ +a_{km}f_{im}+d_{k}u_{ik})
	\end{array}
\end{equation}

これは代数の計算としてやっていくと，非常に煩雑で間違いが起きやすそうです。そこで，少し視覚化してわかりやすくしてみましょう。
多項式の掛け算は，各項目の総当たり戦ですので，列方向に$z_{ij}$，行方向に$z_{ik}$の各要素を置いて，要素同士の組み合わせ表を作ります(図\ref{x03_multiply})。

\begin{figure}[ht]
	\[
		\begin{array}{rl|cccccc|}
			        &                                                       & \multicolumn{6}{c}{z_{ij}=}                                                                                                                                                                                                                                        \\
			        &                                                       & a_{j1}f_{i1} +                                          & a_{j2}f_{i2}+                                            & a_{j3}f_{i3} +                                         & \cdots\cdots +                                          & a_{jm}f_{im} + & d_jU_{ij} \\\hline
			        & \begin{array}{c}a_{k1}f_{i1}\\+\end{array}            & \multicolumn{1}{>{\columncolor{red!9}}c}{}              & \multicolumn{4}{>{\columncolor{cyan!9}}c}{}              & \multicolumn{1}{>{\columncolor{green!9}}c}{}                                                                                                  \\
			z_{ik}= & \begin{array}{c}a_{k2}f_{i2}\\+\end{array}            & \multicolumn{1}{>{\columncolor{cyan!9}}c}{}             & \multicolumn{1}{>{\columncolor{red!9}}c}{\circled{1}}    & \multicolumn{3}{>{\columncolor{cyan!9}}c}{\circled{2}} & \multicolumn{1}{>{\columncolor{green!9}}c}{}                                         \\
			        & \begin{array}{c}a_{k3}f_{i3}\\+\end{array}            & \multicolumn{2}{>{\columncolor{cyan!9}}c}{}             & \multicolumn{1}{>{\columncolor{red!9}}c}{}               & \multicolumn{2}{>{\columncolor{cyan!9}}c}{}            & \multicolumn{1}{>{\columncolor{green!9}}c}{\circled{3}}                              \\
			        & \hspace{10pt}\begin{array}{c}\vdots\\+\end{array}     & \multicolumn{3}{>{\columncolor{cyan!9}}c}{\circled{2}}  & \multicolumn{1}{>{\columncolor{red!9}}c}{}               & \multicolumn{1}{>{\columncolor{cyan!9}}c}{}            & \multicolumn{1}{>{\columncolor{green!9}}c}{}                                         \\
			        & \begin{array}{c}a_{km}f_{im}\\+\end{array}            & \multicolumn{4}{>{\columncolor{cyan!9}}c}{}             & \multicolumn{1}{>{\columncolor{red!9}}c}{}               & \multicolumn{1}{>{\columncolor{green!9}}c}{}                                                                                                  \\
			        & \hspace{7pt}\begin{array}{c}\\d_kU_{ik}\\ \end{array} & \multicolumn{5}{>{\columncolor{green!9}}c}{\circled{3}} & \multicolumn{1}{>{\columncolor{yellow!9}}c}{\circled{4}}                                                                                                                                                 \\\hline
		\end{array}
	\]
	%     \caption{項目同士の総当たりを考える}
	\caption{Consider a round-robin among the items}
	%     \label{x03_multiply}
	Apologies for the misunderstanding, but it seems like you forgot to provide the Japanese text that you want to be translated into English. Please provide the necessary details.
\end{figure}

% 図\ref{x03_multiply}に示されたのは個人$i$についてのものであり，これが人数分ある，すなわち$\sum_{i=1}^N$をつけないといけないことに注意してください。
Please note the diagram \ref{x03_multiply} is about individual $i$, and there are as many as there are people, in other words, you have to put $\sum_{i=1}^N$.
% さて図を軽く色分けしてあるのですが，ここにあるように計算すべき領域を4つに分けて考えていきましょう。
Now, as you can see, I've lightly color-coded the diagram. Let's divide the area we need to calculate into four sections and consider them.
\begin{description}
    \item[\circled{1}の領域] 因子$p$と$p$の積和部分(同じ因子同士の掛け合わせ)
    \item[\circled{2}の領域] 因子$p$と$q$の積和部分(異なる因子同士の掛け合わせ)
    \item[\circled{3}の領域] 因子$p$($q$)と独自因子の積和部分
    \item[\circled{4}の領域] 独自因子同士の積和部分
\end{description}


% この各パートを順に計算していきましょう。まず\circled{1}ですが，たとえば第一因子については
Let's calculate each part in order. First, for \circled{1}, for example, for the first factor,
\begin{equation}
	\frac{1}{N}\sum_{i=1}^{N}a_{j1}a_{k1}F_{i1}^2
\end{equation}

% となることがわかります。ここで，$a_{j1}$と$a_{k1}$は$N$には関係がない($\sum$は$i$が$1$から$N$まで変化することを意味しているが，係数$i$はどちらにも入っていない)ので，総和して割る意味がないことに気づきます。そうなると，必然的にこの式は，
It turns out that $a_{j1}$ and $a_{k1}$ are unrelated to $N$ (the sum means that $i$ varies from $1$ to $N$, but the coefficient $i$ does not enter either), so you'll notice that there's no point in summing and dividing. Then, inevitably, this equation becomes,
\begin{equation}
	\frac{1}{N}\sum_{i=1}^{N}a_{j1}a_{k1}F_{i1}^2=a_{j1}a_{k1}\frac{1}{N}\sum F_{i1}^2
\end{equation}

% となります。この$F_{i1}$は因子得点であり，上の仮定より標準化されたものだということになります。さて，標準得点と標準得点の積和平均は相関係数になることをもう一度思い出してください！そうすると，これは自分自身との相関係数を表していることになりますから，当然$F_{i1}^2=1.0$であることがわかります。
This $F_{i1}$ is a factor score, and from the above assumption, it is standardized. Now, remember that the product mean of standard scores and standard scores becomes a correlation coefficient! So, this represents the correlation coefficient with itself, so naturally it is understood that $F_{i1}^2 = 1.0$.

% 結局，\circled{1}のエリアは
In the end, the area of \circled{1} is
\input{x03_testTheory_escape_07}
% と，とてもあっさり書き下すことができるのです。
And, it can be written very easily.

% つづいて\circled{2}を見てみましょう。ここは異なる因子がかけ合わさった部分ですね。落ち着いて，第一因子と第二因子を例にして考えてみましょう。この箇所で得られるのは，
Let's move on and look at \circled{2}. This is the part where different factors are multiplied. Calmly, let's think about it, taking the first factor and the second factor as examples. What can be achieved at this place is...
\begin{equation}
    \frac{1}{N}\sum a_{j1}F_{i1} a_{k2}F_{i2} = a_{j1}a_{k2}\frac{1}{N}\sum F_{i1}F_{i2}
\end{equation}

% ということになります。ここで，$F_{i1}$および$F_{i2}$はそれぞれ第一，第二因子における個人$i$の因子得点を意味しています。因子得点は標準化されていることをもう一度思い出すと，これは第一因子と第二因子の相関係数になります。ここで，この因子分析が直交因子モデルだと考えますと，因子同士に相関がないわけですから，数字としては$0.0$で消えてしまいます。するとこの部分は，
It comes to this point. Here, $F_{i1}$ and $F_{i2}$ represent the factor scores of individual $i$ in the first and second factors, respectively. Remembering once again that the factor scores are standardized, this becomes the correlation coefficient between the first factor and the second factor. If we assume that this factor analysis is an orthogonal factor model, the factors are uncorrelated, so this figure disappears as $0.0$. Then this part is,

\input{x03_testTheory_escape_09}
% となることがわかりました。つまり，この領域\circled{2}は，すべて$0$になってしまうのです。
It was found out. In other words, this area marked as \circled{2} all ends up being $0$.


% 続いて\circled{3}の部分について考えてみましょう。これはある共通因子と独自因子の積和部分です。例によって標準得点同士の関係から，相関係数を算出することになりますが，独自因子は共通因子と無相関であることを考えると，
Let's consider the part \circled{3} next. This is the sum of products of a common factor and a unique factor. As with the example, we'll calculate the correlation coefficient from the relationship between standard scores, but considering that the unique factor is uncorrelated with the common factor,
\begin{equation}
	\frac{1}{N}\sum a_{i1}F_{i1}d_{j}U_{ij}=a_{ik}d_{j} \frac{1}{N}\sum U_{ij}F_{i1} =0
\end{equation}

% とこのように，この箇所もすべて$0$になってしまいます。
And so, all of this part also becomes $0$.

% 最後の\circled{4}に至っては，独自因子と独自因子の積和ですから，これも
As for the final \circled{4}, since it is the product sum of unique factors, this too
\begin{equation}
    \frac{1}{N}\sum d_{j}d_{k}U_{ij}U_{ik}=d_{j}d_{k}\frac{1}{N}\sum U_{ij}U_{ik}=0
\end{equation}

% のように$0$になります。結局，消えて無くなるのがほとんどで，残るのは\textcircled{1}の部分だけであり，$r_{jk}$を考えるときはそこだけ考慮すればよいことになります。
It becomes $0$ like this. After all, most of them disappear, and the only part left is the \textcircled{1} part, which means that we only need to consider that part when considering $r_{jk}$.

% 整理すると，
To summarize,
\begin{equation}
    r_{jk}=a_{j1}a_{k1}+a_{j2}a_{k2}+ \cdots +a_{jm}a_{km}
\end{equation}

% ということがわかります。つまり，項目$j$と項目$k$の相関係数は，項目$j$の因子負荷量と項目$k$の因子負荷量を，すべての因子について総和したものであるということです。因子分析の基本モデルから導出されるこの定理を，とくに\textbf{因子分析の第二定理}と呼びます。
It is understood that the correlation coefficient of Item j and Item k is the sum of the factor loadings of Item j and Item k for all factors. This theorem derived from the basic model of factor analysis is particularly called the \textbf{second theorem of factor analysis}.

% ここで同じ項目同士の相関を考えてみましょう。項目$j$と項目$j$の相関係数は，もちろん$1.0$になりますね。これを因子分析の基本式で表すと，次のように表現できます。
Let's think about the correlation between the same items here. The correlation coefficient between item $j$ and item $j$, of course, will be $1.0$. This can be expressed in the basic equation of factor analysis as follows.
\begin{equation}
	\label{facteo1}
	r_{jj}=a_{j1}^2+a_{j2}^2+ \cdots + a_{jm}^2 + d_j^2 = 1.0
\end{equation}


% さて，この式が意味するのはなんでしょうか。意味を考えてみると，ある項目それ自身との相関係数は，因子負荷の二乗和からなっている，ということがわかります。これこそ\textbf{因子分析の第一定理}と呼ばれるものであり，解けるはずのなかった方程式を解くための鍵となる式なのです。
So, what does this equation mean? When we think about its meaning, we can understand that the correlation coefficient of an item with itself comes from the sum of the squares of factor loadings. This is called the \textbf{First Theorem of Factor Analysis}, and it is the key equation for solving a supposed unsolvable equation.

% \subsection{因子分析の定理}
\subsection{Theorem of Factor Analysis}
% 数式の展開はいったんここまでにして，
Let's put a pause on the expansion of the formula for now,
% 第一定理は次のような形をしているのでした。
The first theorem was in the following form.
\[a_{j1}^2+a_{j2}^2+ \cdots + a_{jm}^2 + d_j^2 = 1.0 \]
ここで共通因子部分を，$a_{j1}^2+a_{j2}^2+ \cdots + a_{jm}^2=h_j^2$のようにすると，この式は単に$h_j^2 + d_j^2 = 1.0$となります。この$h_j^2$のことをとくに\keyterm{共通性}{きょうつうせい}{communality}といいますが，この式は共通性と独自因子の二乗和が$1.0$になることを意味しています。言い換えると，全体を100\%とした比率で共通性と誤差を比較できるということです。
共通性は因子負荷量の二乗和で，共通因子はそのテストの背後にある共通の要因，すなわちテストで測定したいものだったわけです。古典的テスト理論では，モデル式の展開から信頼性を全分散中に示る真のスコアの割合と定義しましたが，因子分析モデルはこのように1つの項目$j$における共通因子の割合を算出し，項目の信頼性を考えることができます。因子分析モデルにおける信頼性は，1項目の中に含まれる共通因子の大きさだとも言えるわけです。逆に$d_j^2=1-h_j^2$は\keyterm{独自性}{どくじせい}{uniqueness}は，当該項目がそのテストで測っていないものの大きさを表しており，この割合があまりにも大きいと「この項目は全然関係ないものを測っちゃってるんじゃないか？ 」と疑われることになります。多因子モデルにおいては，多角的に対象を切り分けるために多くの質問を投げかけるわけですが，独自性の高い項目は回答者に負担をかけるだけの邪魔なものですから，実践上はこうした項目を除外することが少なくありません。因子分析には\keyterm{単純構造の原則}{たんじゅんこうぞうのげんそく}{principle of simple structure}と呼ばれるものがあり，項目は該当する因子を適切に反映し，かつ，他の因子と関係ないことが美しいとされます。尺度構成段階では，共通性(独自性)をみて項目の良し悪しが判断されるのです。

次に第二定理を見てみましょう。第二定理は次のような形をしているのでした。
\[  r_{jk}=a_{j1}a_{k1}+a_{j2}a_{k2}+ \cdots +a_{jm}a_{km} \]
% 2つの項目$j$と$k$の相関係数は，それぞれの因子負荷量の積和の形で表される，というものです。ここに誤差の話は入ってこず，共通因子だけで話ができています。
The correlation coefficient between two items $j$ and $k$ is represented as the sum of the product of their respective factor loadings. In this context, there is no mention of error, and the discussion can be held with only common factors.

% 左辺の相関係数は，2つの項目がどれほど同じものを測定しているかの指標です。相関係数(の絶対値)が高ければ高いほど，2つの項目は同じものを指し示しているわけです。逆に相関係数が低いということは，2つの項目に関係がないことを表します。ここで右辺に目をやりますと，右辺の各項目は因子負荷量の積の形になっています。左辺の値が小さくなる1つの理由は，ある共通因子$m$が項目$j,k$に対して，異なる方向で寄与しているからだと考えられるでしょう。そしてそのパターンが一貫していないという状況です。そもそも相関係数が小さいところからは因子を見つけ出すのは難しいのですが，そうした状況があるのはある項目ペアについて因子同士の向きがバラバラに影響しているからだと言えます。そのような状況は，測定がきちんとできているかどうか怪しいですね。\textbf{測定の一義性}とも言われますが
The correlation coefficient on the left side is an indicator of how much the two items are measuring the same thing. The higher the (absolute value of the) correlation coefficient, the more the two items indicate the same thing. Conversely, a low correlation coefficient indicates that there is no relationship between the two items. If you look at the right side, each item on the right side is in the form of a product of factor loadings. One reason for the value on the left to be small can be considered because a common factor $m$ contributes differently to items $j,k$. And it's a situation where that pattern is not consistent. It's difficult to find factors from where the correlation coefficient is small in the first place, but such a situation exists because the directions of the factors are affecting variously about some item pair. Such circumstances are suspicious as to whether the measurement is being done properly. It's also said to be \textbf{uniqueness of measurement}.
% \footnote{公理論的測定論における測定論の基礎問題には，表現の問題，一意性の問題，有意味性の問題，尺度構成の問題があり\parencite{Coombs1970}，一意性の問題とは「測定手順があるとして，対象に数値を割り当てる時にどの程度まで自由に行えるか」あるいは「与えた数値の相互変換はどのようなものが許容されるか」といったことを指します。}
% ，そのような尺度は妥当性が低いといえるでしょう。
It could be said that such a measure has low validity.

% このように，因子分析モデルは第一定理で信頼性を，第二定理で妥当性をあらわすものになっているのです。
In this way, the factor analysis model represents reliability in the first theorem and validity in the second theorem.

% \section{課題}
\section{Tasks}
% \paragraph{テスト理論と因子分析} 因子分析モデルは古典的テスト理論をどのように発展させたのか，添字に注意しながら数式で表現してみよう。
\paragraph{Test Theory and Factor Analysis} Let's try to express in mathematical formulas, with care to indices, how the factor analysis model advances the classical test theory.
% \paragraph{因子分析の定理の導出} 因子分析の定理の導出を自分でできるようになろう。
\paragraph{Derivation of the Theorem of Factor Analysis} Let's learn to derive the theorem of factor analysis by ourselves.
% \paragraph{因子分析の定理の意味} 因子分析の定理が何を意味しているのか，自分なりの言葉で説明してみよう。
\paragraph{The Meaning of The Factor Analysis Theorem} Let me try to explain in my own words what the Factor Analysis Theorem means.

\end{document}
